% Source: https://artofproblemsolving.com/community/c3h3267915_2023_lmt_fall_guts_round_p16p27___lexington_mathematical_tournament
% Stable identity: aops-topic-3267915-q22
Take $B=(0,0),C=(14,0),A=(9,12)$; the given distances are verified by the Pythagorean theorem. The signed side-length formulas for excentres give
\[I_A=\dfrac{-14A+13B+15C}{-14+13+15}=(6,-12),\quad I_B=\dfrac{14A-13B+15C}{14-13+15}=(21,\dfrac{21}{2}),\]
\[I_C=\dfrac{14A+13B-15C}{14+13-15}=(-7,14).\]
These formulas follow from the internal and external angle-bisector ratios (equivalently, equal distances to the three side lines with the indicated signs). The triangle area is $84$. The excentral area, using vectors $(15,\dfrac{45}{2})$ and $(-13,26)$ from $I_A$, is
\[\dfrac{1}{2}\left(15(26)+13\cdot\dfrac{45}{2}\right)=\dfrac{1365}{4}.\]
The ratio is $\boxed{\dfrac{65}{16}}$.
\begin{asy}
usepackage("amsmath"); size(220);
pair A=(9,12),B=(0,0),C=(14,0),IA=(6,-12),IB=(21,10.5),IC=(-7,14);draw(A--B--C--cycle,linewidth(1));draw(IA--IB--IC--cycle,blue);draw(A--IA,dashed);draw(B--IB,dashed);draw(C--IC,dashed);label("$A$",A,N);label("$B$",B,SW);label("$C$",C,SE);label("$I_A$",IA,S);label("$I_B$",IB,E);label("$I_C$",IC,W);label("$15$",(A+B)/2,NW);label("$13$",(A+C)/2,NE);label("$14$",(B+C)/2,S);
\end{asy}
